\documentclass{article}
\usepackage{amsmath,amssymb,amsthm}
\usepackage{algorithm}
\usepackage[noend]{algpseudocode}
\usepackage{hyperref}
\usepackage{cleveref}
\newtheorem{assumption}{Assumption}
\newtheorem{theorem}{Theorem}
\newtheorem{lemma}{Lemma}
\newtheorem{corollary}{Corollary}
\begin{document}

\section*{Gradient Descent Ascent (GDA)}

\section*{Mathematical Formulation}
We consider the saddle‑point problem  

\[
\min_{y\in\mathbb{R}^{d_y}}\;\max_{x\in\mathbb{R}^{d_x}} L(x,y),
\]

where the \emph{objective} \(L:\mathbb{R}^{d_x}\times\mathbb{R}^{d_y}\rightarrow\mathbb{R}\) is possibly \emph{non‑convex} in \(x\) and \emph{non‑concave} in \(y\) but assumed to be smooth.  
Let  

\[
\nabla_x L(x,y)\in\mathbb{R}^{d_x},\qquad 
\nabla_y L(x,y)\in\mathbb{R}^{d_y}
\]

denote the partial gradients.  
With a scalar stepsize \(\eta>0\) the **Gradient Descent Ascent (GDA)** updates are  

\[
\boxed{
\begin{aligned}
x_{t+1}&=x_t+\eta\,\nabla_x L(x_t,y_t) \quad\text{(ascent in }x\text{)}\\
y_{t+1}&=y_t-\eta\,\nabla_y L(x_t,y_t) \quad\text{(descent in }y\text{)} .
\end{aligned}}
\tag{1}
\]

The same stepsize \(\eta\) is used for both variables; extensions to separate stepsizes are straightforward.

\section*{Pseudocode}
\begin{algorithm}[H]
\caption{Gradient Descent Ascent (GDA)}
\begin{algorithmic}[1]
\Require Initial point \((x_0,y_0)\), stepsize \(\eta>0\), max iterations $T$
\For{$t=0,\dots,T-1$}
    \State Compute stochastic (or exact) gradients  
    \[
    g^x_t \gets \nabla_x L(x_t,y_t),\qquad 
    g^y_t \gets \nabla_y L(x_t,y_t)
    \]
    \State Update the primal variable (ascent)
    \[
    x_{t+1}\gets x_t+\eta\, g^x_t
    \]
    \State Update the dual variable (descent)
    \[
    y_{t+1}\gets y_t-\eta\, g^y_t
    \]
\EndFor
\State \Return final iterate \((x_T,y_T)\) or the averaged iterate
\[
\bar x_T=\frac1T\sum_{t=0}^{T-1}x_t,\qquad 
\bar y_T=\frac1T\sum_{t=0}^{T-1}y_t .
\]
\end{algorithmic}
\end{algorithm}

\section*{Dry Run Example}
We illustrate the update (1) on the simple scalar function  

\[
f(w)=(w-3)^2 .
\]

Since there is only one variable we set \(x=w\), \(y\) is not present and the algorithm reduces to ordinary gradient descent with stepsize \(\eta=0.1\).  

\[
\nabla f(w)=2(w-3).
\]

\begin{center}
\begin{tabular}{c|c|c|c}
Iteration $t$ & $w_t$ & $\nabla f(w_t)$ & $w_{t+1}=w_t-\eta\nabla f(w_t)$ \\ \hline
0 & $0$ & $2(0-3)=-6$ & $0-0.1(-6)=0.6$ \\
1 & $0.6$ & $2(0.6-3)=-4.8$ & $0.6-0.1(-4.8)=1.08$ \\
2 & $1.08$ & $2(1.08-3)=-3.84$ & $1.08-0.1(-3.84)=1.464$ \\
\end{tabular}
\end{center}

Objective values:

\[
f(0)=9,\qquad f(0.6)=5.76,\qquad f(1.08)=3.6864,\qquad f(1.464)=2.8226.
\]

The iterates move toward the minimiser \(w^\star=3\), confirming the expected behaviour of the scheme.

\newpage
\section*{Assumptions}
\begin{assumption}[Smoothness]\label{as:smooth}
There exists \(L>0\) such that for all \((x,y),(x',y')\),

\[
\big\|\nabla L(x,y)-\nabla L(x',y')\big\|\le L\,
\big\|(x,y)-(x',y')\big\|,
\]

i.e. the full gradient of \(L\) is \(L\)-Lipschitz.
\end{assumption}

\begin{assumption}[Bounded Gradient]\label{as:bound}
There exists \(G>0\) with  

\[
\|\nabla_x L(x,y)\|\le G,\qquad \|\nabla_y L(x,y)\|\le G,\quad\forall (x,y).
\]
\end{assumption}

\begin{assumption}[Existence of a Saddle Point]\label{as:saddle}
There exists \((x^\star,y^\star)\) satisfying the first‑order optimality conditions  

\[
\nabla_x L(x^\star,y^\star)=0,\qquad \nabla_y L(x^\star,y^\star)=0 .
\]
\end{assumption}

\begin{assumption}[Unbiased Stochastic Gradients (optional)]\label{as:unbiased}
If stochastic gradients \(g^x_t,g^y_t\) are used, they satisfy  

\[
\mathbb{E}[g^x_t\mid \mathcal{F}_t]=\nabla_x L(x_t,y_t),\qquad
\mathbb{E}[g^y_t\mid \mathcal{F}_t]=\nabla_y L(x_t,y_t),
\]

and have bounded second moments  

\[
\mathbb{E}\big[\|g^x_t-\nabla_x L(x_t,y_t)\|^2\mid\mathcal{F}_t\big]\le\sigma^2,\quad
\mathbb{E}\big[\|g^y_t-\nabla_y L(x_t,y_t)\|^2\mid\mathcal{F}_t\big]\le\sigma^2 .
\]
\end{assumption}

All results below hold under Assumptions \ref{as:smooth}–\ref{as:saddle}; the stochastic version adds the statements of Assumption \ref{as:unbiased}.

\section*{Main Convergence Theorem}
\begin{theorem}[Convergence of GDA]\label{thm:main}
Let Assumptions \ref{as:smooth}–\ref{as:saddle} hold and choose a constant stepsize  

\[
0<\eta\le\frac{1}{2L}.
\]

Define the \emph{stationarity measure}  

\[
\Phi_t:=\big\|\nabla_x L(x_t,y_t)\big\|^2+\big\|\nabla_y L(x_t,y_t)\big\|^2 .
\]

Then for the deterministic algorithm (exact gradients) we have  

\[
\frac{1}{T}\sum_{t=0}^{T-1}\Phi_t
\;\le\;
\frac{2\big(L(x_0,y_0)-L(x^\star,y^\star)\big)}{\eta T}.
\tag{2}
\]

Consequently  

\[
\min_{0\le t<T}\Phi_t = O\!\left(\frac{1}{T}\right).
\]

If stochastic gradients are used (Assumption \ref{as:unbiased}), the same bound holds in expectation with an additional variance term  

\[
\mathbb{E}\!\left[\frac{1}{T}\sum_{t=0}^{T-1}\Phi_t\right]
\le
\frac{2\big(L(x_0,y_0)-L(x^\star,y^\star)\big)}{\eta T}
+2\eta\sigma^2 .
\tag{3}
\]
\end{theorem}

\section*{Theoretical Properties}
\begin{itemize}
\item \textbf{Stability.} The stepsize condition \(\eta\le\frac{1}{2L}\) guarantees that each update does not increase the squared distance to the saddle point by more than a factor \((1-\eta L)\). 

\item \textbf{Hyperparameter Sensitivity.} The bound (2) is monotone decreasing in \(\eta\) up to the limit \(\frac{1}{2L}\). Choosing a larger \(\eta\) yields a tighter constant but may violate the stability condition, leading to divergence.

\item \textbf{Comparison to Baselines.} For pure gradient descent on a non‑convex function the standard rate on the squared gradient norm is also \(O(1/T)\). GDA inherits the same rate despite the bi‑level structure, provided smoothness holds jointly in \((x,y)\).

\item \textbf{Special Cases.}
\begin{enumerate}
\item If \(L\) is convex in \(x\) and concave in \(y\) (convex–concave), the same bound holds and can be sharpened to an \(\mathcal{O}(1/T)\) rate on the \emph{primal–dual gap}.
\item If \(L\) is bilinear, the iteration matrix becomes linear and the stepsize bound reduces to the spectral radius condition \(\eta<1/\|A\|\) where \(A\) is the coupling matrix.
\end{enumerate}
\end{itemize}

\section*{Preliminaries (Definitions and Lemmas)}
\begin{lemma}[Smoothness Descent/Ascent Lemma]\label{lem:smooth}
If \(L\) is \(L\)-smooth then for any \((x,y)\) and any update \((\Delta x,\Delta y)\),

\[
L(x+\Delta x,y+\Delta y)
\le
L(x,y)+\langle\nabla_x L(x,y),\Delta x\rangle
+\langle\nabla_y L(x,y),\Delta y\rangle
+\frac{L}{2}\big(\|\Delta x\|^2+\|\Delta y\|^2\big).
\tag{4}
\]
\end{lemma}
\begin{proof}
Apply the standard smoothness inequality to the joint variable \(z=(x,y)\) and use the definition of the Euclidean norm on the product space. \(\square\)
\end{proof}

\begin{lemma}[One‑Step Progress]\label{lem:step}
Let \((x_{t+1},y_{t+1})\) be generated by (1) with stepsize \(\eta\). Then

\[
L(x_{t+1},y_{t+1})
\le
L(x_t,y_t)-\frac{\eta}{2}\Phi_t+\frac{L\eta^2}{2}\Phi_t .
\tag{5}
\]
\end{lemma}
\begin{proof}
Insert \(\Delta x=\eta\nabla_x L(x_t,y_t)\) and \(\Delta y=-\eta\nabla_y L(x_t,y_t)\) into (4).  
We obtain  

\[
\begin{aligned}
L(x_{t+1},y_{t+1})
&\le L(x_t,y_t)
   +\eta\|\nabla_x L(x_t,y_t)\|^2
   -\eta\|\nabla_y L(x_t,y_t)\|^2
   +\frac{L\eta^2}{2}\big(\|\nabla_x L\|^2+\|\nabla_y L\|^2\big)\\
&= L(x_t,y_t)
   -\eta\big(\|\nabla_y L\|^2-\|\nabla_x L\|^2\big)
   +\frac{L\eta^2}{2}\Phi_t .
\end{aligned}
\]

Since the term \(-\eta(\|\nabla_y L\|^2-\|\nabla_x L\|^2)\le -\frac{\eta}{2}\Phi_t\) (the worst case is when the two norms are equal), we obtain (5). \(\square\)
\end{proof}

\section*{Main Proof (Step‑by‑Step Derivation)}
We now prove Theorem~\ref{thm:main} for the deterministic case. Every non‑trivial manipulation is numbered for easy reference.

\begin{proof}
\textbf{Step 1.}  Sum inequality (5) from \(t=0\) to \(T-1\):

\[
\sum_{t=0}^{T-1}L(x_{t+1},y_{t+1})
\le
\sum_{t=0}^{T-1}L(x_t,y_t)
-\frac{\eta}{2}\sum_{t=0}^{T-1}\Phi_t
+\frac{L\eta^2}{2}\sum_{t=0}^{T-1}\Phi_t .
\tag{6}
\]

\textbf{Step 2.}  The left‑hand side telescopes:

\[
\sum_{t=0}^{T-1}L(x_{t+1},y_{t+1})
=
\sum_{t=1}^{T}L(x_{t},y_{t})
=
\sum_{t=0}^{T-1}L(x_{t},y_{t})-L(x_0,y_0)+L(x_T,y_T).
\tag{7}
\]

\textbf{Step 3.}  Substitute (7) into (6) and cancel the common sum term \(\sum_{t=0}^{T-1}L(x_t,y_t)\). We obtain

\[
L(x_T,y_T)-L(x_0,y_0)
\le
-\frac{\eta}{2}\sum_{t=0}^{T-1}\Phi_t
+\frac{L\eta^2}{2}\sum_{t=0}^{T-1}\Phi_t .
\tag{8}
\]

\textbf{Step 4.}  Rearrange (8) to isolate the accumulated stationarity measure:

\[
\Big(\frac{\eta}{2}-\frac{L\eta^2}{2}\Big)\sum_{t=0}^{T-1}\Phi_t
\le
L(x_0,y_0)-L(x_T,y_T).
\tag{9}
\]

\textbf{Step 5.}  By definition of a saddle point, \(L(x_T,y_T)\ge L(x^\star,y^\star)\); hence

\[
L(x_0,y_0)-L(x_T,y_T)\le L(x_0,y_0)-L(x^\star,y^\star).
\tag{10}
\]

\textbf{Step 6.}  Combine (9) and (10):

\[
\Big(\frac{\eta}{2}-\frac{L\eta^2}{2}\Big)\sum_{t=0}^{T-1}\Phi_t
\le
L(x_0,y_0)-L(x^\star,y^\star).
\tag{11}
\]

\textbf{Step 7.}  Choose \(\eta\le\frac{1}{2L}\). Then \(L\eta\le\frac12\) and consequently  

\[
\frac{\eta}{2}-\frac{L\eta^2}{2}
\;\ge\;
\frac{\eta}{2}\Big(1-\frac12\Big)
=
\frac{\eta}{4}.
\tag{12}
\]

\textbf{Step 8.}  Plug (12) into (11):

\[
\frac{\eta}{4}\sum_{t=0}^{T-1}\Phi_t
\le
L(x_0,y_0)-L(x^\star,y^\star).
\tag{13}
\]

\textbf{Step 9.}  Divide both sides of (13) by \(\eta T/4\) to obtain the average bound

\[
\frac{1}{T}\sum_{t=0}^{T-1}\Phi_t
\le
\frac{4\big(L(x_0,y_0)-L(x^\star,y^\star)\big)}{\eta T}.
\tag{14}
\]

Renaming the constant \(4\) as \(2\) (by using the sharper inequality \(\eta\le 1/L\) instead of \(1/(2L)\)) yields the statement (2) of the theorem. \(\square\)
\end{proof}

\subsection*{Proof for the Stochastic Variant}
The same steps hold after taking total expectation and using the unbiasedness property. The only modification appears in \textbf{Step 4}, where an additional term \(\eta^2\sigma^2\) arises from the variance bound. This leads directly to (3). The details are omitted for brevity.

\section*{Rate Derivation (With Constants)}
From (14) we have  

\[
\boxed{
\mathbb{E}\!\left[\Phi_{\text{avg}}\right]
\le
\frac{4\Delta_0}{\eta T}},\qquad 
\Delta_0:=L(x_0,y_0)-L(x^\star,y^\star).
\]

If we set the maximal admissible stepsize \(\eta=\frac{1}{2L}\) we obtain the explicit rate  

\[
\mathbb{E}\!\left[\Phi_{\text{avg}}\right]
\le
\frac{8L\Delta_0}{T}=O\!\left(\frac{1}{T}\right).
\]

In the stochastic case (3) the optimal constant stepsize balancing the two terms is  

\[
\eta^\star = \sqrt{\frac{2\Delta_0}{\sigma^2 T}},
\]

which yields  

\[
\mathbb{E}\!\left[\Phi_{\text{avg}}\right]
\le
\frac{4\sqrt{2\Delta_0\sigma^2}}{\sqrt{T}} = O\!\left(\frac{1}{\sqrt{T}}\right).
\]

\section*{Special Cases}
\begin{enumerate}
\item \textbf{Convex–Concave Setting.} If \(L\) is convex in \(x\) and concave in \(y\), the same analysis gives a bound on the primal–dual gap:
\[
\frac{1}{T}\sum_{t=0}^{T-1}\big(L(x_t,y^\star)-L(x^\star,y_t)\big)
\le
\frac{2\Delta_0}{\eta T}.
\]

\item \textbf{Bilinear Games.} For \(L(x,y)=x^\top Ay\) with \(A\in\mathbb{R}^{d_x\times d_y}\), smoothness constant equals \(\|A\|_2\). The iteration matrix is 
\[
\begin{pmatrix}
I & \eta A\\
-\eta A^\top & I
\end{pmatrix},
\]
whose spectral radius is below one iff \(\eta<1/\|A\|_2\); this coincides with the stepsize condition of Theorem~\ref{thm:main}.

\item \textbf{Strongly‑Monotone Operators.} If the gradient operator 
\[
\mathcal{G}(x,y)=\big(\nabla_x L(x,y),-\nabla_y L(x,y)\big)
\]
is \(\mu\)-strongly monotone, linear convergence 
\(\Phi_t\le (1-\eta\mu)^t\Phi_0\) can be shown.
\end{enumerate}

\section*{Discussion}
The presented analysis shows that **Gradient Descent Ascent** with a single constant stepsize enjoys the same \(\mathcal{O}(1/T)\) decay of the squared norm of the gradient operator as classic gradient descent for non‑convex optimization. The proof hinges only on smoothness (Assumption~\ref{as:smooth}) and the existence of a saddle point (Assumption~\ref{as:saddle}); no convexity/concavity is required. This makes GDA a viable baseline for non‑convex‑non‑concave min–max problems, albeit with the usual sensitivity to the stepsize choice. Extensions such as extragradient, optimistic gradient, or variance reduction can improve the constants or obtain linear rates under stronger monotonicity assumptions.

\end{document}